\SourceLeaf{131}
Wo man übrigens mit sieben oder zehn Decimalen ausreicht, gewähren freilich die gewöhnlichen oder die \textsc{Vlacq}schen, auch von \textsc{Vega} herausgegebenen Tafeln bei weitem mehr Bequemlichkeit, als die \textsc{Leonelli}schen. Und in der That sind jene auch bei den allerfeinsten astronomischen oder sonst auf die Körperwelt sich beziehenden Rechnungen, ohne Ausnahme, überflüssig genau. Bei analytischen Rechnungen, oder auch bei der primitiven Construction von Tafeln, kommt man jedoch öfters in den Fall, wo man eine grössere Schärfe wünscht, und dann, wenn man anders nicht noch mehr als 13 oder 18 Ziffern verlangt, sind \textsc{Leonellis} Tafeln unstreitig das Brauchbarste und Bequemste, was man zu diesem Behufe anwenden kann.

Dem \textsc{Delambre}schen Berichte zufolge hat übrigens \textsc{Brigg} ganz dieselbe Methode in seiner \emph{arithmetica logarithmica} vorgetragen, und eine ähnliche Hülfstafel auf 14 Decimalen berechnet; allein die französischen Geometer selbst haben von diesem seltenen Werke nur Ein Exemplar auftreiben können, worin gerade dieser Theil gefehlt hat. Das Verdienst der eigenen Erfindung und eigenen Berechnung bleibt also \textsc{Leonelli} immer ungeschmälert, und es ist billig, dass wir ihm besonders für letztere den gebührenden Dank zollen.

Den Zweck der zweiten Abhandlung haben wir bereits oben angezeigt; wir wollen nun sehen, wie \textsc{Leonelli} denselben erreichen will. Er schlägt eine aus drei Columnen bestehende Tafel vor; um uns kürzer zu fassen, wollen wir, anstatt \textsc{Leonellis} schwerfällige und unnöthige Terminologie zu gebrauchen, die zusammengehörigen Glieder dieser drei Columnen durch $P$, $Q$, $R$ bezeichnen. Diese Grössen sollen so von einander abhängen, dass, wenn man die Zahl, deren Logarithm $P$ ist, durch $x$ bezeichnet,
\[
Q = \log\left(1+\frac{1}{x}\right), \qquad R = \log(1+x)
\]
sei; daher immer $R = P+Q$ sein wird; ferner soll diese Tafel nicht nach $x$, sondern nach $P$ geordnet sein, oder $P$ soll gleichförmig wachsen, und zwar von Null bis ins Unendliche, oder vielmehr bis $Q$ als verschwindend betrachtet werden kann. Der Gebrauch einer solchen Tafel lässt sich leicht übersehen. Soll aus $\log a$ und $\log b$ der Logarithm von $a+b$ bestimmt werden: so geht man (wenn man voraussetzt, dass $a$ grösser ist als $b$) mit $\log a - \log b$ in die erste Columne ein, oder setzt diese Differenz $= P$: dann ist offenbar
\[
\log(a+b) = Q + \log a = R + \log b.
\]

